\documentclass[twocolumn, 10pt]{jsarticle}
% upLaTeX/dvipdfmx環境
\usepackage[top=15truemm,bottom=15truemm,left=10truemm,right=10truemm]{geometry}
\pagestyle{empty}
\usepackage[dvipdfmx]{graphicx}
\usepackage[dvipdfmx]{color}
\usepackage[dvipdfmx]{xcolor}
\usepackage{amsmath}
\usepackage{amssymb}

\title{付録：生成AIの活用方法}
\author{23266025}
\date{2026年7月31日}

\begin{document}
\maketitle

\section*{使用したツール}

Claude Code（Anthropic社の生成AIモデル Claude Sonnet 5 を搭載したCLIエージェント）を使用した．対話を通じて，関連研究の調査・PDFの抽出・原文との照合などをAIに実行させながら，要約を作成した．

\section*{活用方法}

\subsection*{1. 対象論文の選定}

課題の対象論文誌・トピック・期間（2016年以降）のもとで，筆者が最近読んだ論文から挙げた3名の研究者（Danilo P. Mandic，Mahnaz Arvaneh，Fabien Lotte）の業績から候補を探すよう指示した．Claude Codeはdblp・Google Scholar等で3名を独立に調査し，候補をリスト化した．調査結果をもとに，Mandic氏の論文を選定した．

\subsection*{2. PDFの抽出と原文照合}

選定した論文のIEEE出版版PDFから，\texttt{pdftotext}ベースのスクリプトで読み順テキストとレイアウト確認用テキストを抽出させた．抽出結果は下書きにすぎないため，論文中の図（Fig.\ 1〜3）も画像として読み取らせ，整合性を確認した．

\subsection*{3. プレプリント版と出版版の突き合わせ}

IEEE出版版PDFを入手した際，arXivプレプリント版（v1）と突き合わせ，出版前後の変化を調べさせた．その結果，（a）応用例が3例から2例に差し替え，（b）局所最適解が大域最適解でもあるとする定理を削除，（c）四元数ニュートン法の2次収束性に関する定理を追加，（d）数値実験（QTFMによる収束検証）を追加，の4点を検出させた．査読前後の変化を把握したことで，著者らが最終的に主張したかった内容を正確に理解できた．

\subsection*{4. 本文との対応位置の検証}

要約中の各主張について，原論文中で根拠となるSection・Theorem・Proposition・式番号・Figureを逐一確認し，本文中に付記した．有効性の検証（第4〜5章）で用いた数値（テスト正解率約96\%など）は，Fig.\ 2・Fig.\ 3の画像を直接読み取らせて確認したうえで記述した．

\subsection*{5. 引用文献の系譜調査}

原論文が先行研究として挙げているQi ら（2022）・Flamant ら（2022）・Liu ら（2020）については，各論文自体は入手せず，原論文自身のIntroduction節にある内容紹介（Qi らの1階・2階微分の計算方法，Flamant らとLiu らによる1階の特徴づけなど）を根拠に，各研究の位置づけをAIに整理させた．一方，原論文が基盤とするGHR微分の枠組み（Xu ら 2015, \textit{Royal Society Open Science}；Xu, Xia, and Mandic, 2016, \textit{IEEE TNNLS}）については，参考文献リストの著者欄にDongpo Xuの名前が含まれることをAIに確認させ，責任著者自身による先行研究であることを特定させた．

\subsection*{6. 文章の構成・文体の反復修正}

日本語の学術文章として明確かつ簡潔になるよう，章立て・文体を複数回にわたり見直させた．章立てを論文自身の構成（背景\(\to\)理論\(\to\)応用\(\to\)結論）に合わせる，抽象名詞句ではなく能動態で記述する，間接疑問文を名詞句に言い換える，といった修正を都度指示しながら適用させた．

\section*{生成AIによる出力と，筆者自身が行ったこと}

要約の観点や構成の設計を含め，以下の判断はすべて筆者が下した．

\begin{itemize}
    \item 要約の観点（GHR微分の枠組みとその適用先）と章構成の概形の決定
    \item 候補論文の中からの最終的な論文選定
    \item 主語と述語の対応関係，修飾語の係り受け，章立ての構成原理など，文章の根本的な構造に関する指摘と修正指示
    \item 執筆者自身が作成した文章表現をAIに整形させ，本文へ組み込ませたこと
    \item 数式・定理・数値の一致確認にとどまらず，論文の展開を追い，各応用が定理の要求条件（二次形式・強凸性など）を満たす点・満たさない点を把握し，読者に伝わるよう記述したこと
\end{itemize}

以上のように，本課題では生成AIを，原論文の抽出・照合・下書き作成・反復的な文章校正のための作業ツールとして活用した．一方で，要約に含める内容の取捨選択と原論文の内容の解釈については，筆者自身が原論文を読み込んだうえで判断した．

\section*{総作業時間}

本課題全体でのAIとの共同作業には，合計で11時間34分を要した．

\end{document}
